\documentclass[documento]{gpt}

% ---------------------------------------------------------------------
% Dati specifici del documento
% ---------------------------------------------------------------------
\titolo{Dichiarazione di Impegni}
\uso{esterno}
\destinatari{Prof.~Tullio Vardanega\newline Prof.~Riccardo Cardin}

% ---------------------------------------------------------------------
% Registro delle modifiche
% ---------------------------------------------------------------------
% Sintassi: \modifica{versione}{AAAA-MM-GG}{autore}{verificatore}{descrizione}
% La riga più recente va messa in cima.
\modifica{1.1.0}{2026-10-27}{spagnoli}{corro}{modifica data completamento}
\modifica{1.0.0}{2026-10-27}{--}{corro}{Verifica}
\modifica{0.1.0}{2026-10-27}{spagnoli, cescobolla}{corro}{Prima stesura del documento}

\begin{document}

% ---------------------------------------------------------------------
% Contenuto del Documento
% ---------------------------------------------------------------------
\section{Candidatura}

\subsection{Capitolato Selezionato}
A seguito dell'analisi preliminare e del colloquio conoscitivo svolto con l'azienda proponente Eggon S.r.l., il gruppo ha individuato nel capitolato C2 -- NEXUM Concierge la scelta ottimale per lo svolgimento del progetto didattico. La decisione si basa sui seguenti fattori determinanti:

\begin{itemize}
    \item \textbf{Chiarezza della documentazione e perimetrazione dell'MVP:} La documentazione fornita dall'azienda si distingue per la precisione nella definizione dei requisiti funzionali minimi (gestione di ferie/permessi e consultazione di regole e manuali aziendali). Questa chiara delimitazione riduce significativamente i rischi di incertezza e di espansione incontrollata dei requisiti.
    
    \item \textbf{Disponibilità e supporto dell'azienda proponente:} Il colloquio conoscitivo con il referente aziendale Gianpaolo ha confermato una spiccata disponibilità al dialogo costruttivo e al supporto continuo, elemento fondamentale per la risoluzione tempestiva di eventuali dubbi architetturali ed operativi durante lo sviluppo.
\end{itemize}

\subsection{Assunzione Impegni}
Con la presente, il Git Push Team comunica formalmente la propria candidatura per l'aggiudicazione del capitolato C2 -- NEXUM Concierge, proposto da Eggon S.r.l.

A fronte di tale candidatura, il gruppo Git Push Team si impegna formalmente a:
\begin{itemize}
    \item Completare il progetto didattico entro il 23/03/2027
    \item Rispettare rigorosamente il preventivo economico presentato
    \item Garantire il monte ore di lavoro dichiarato da ciascun membro del team
    \item Seguire la metodologia di lavoro (\textit{way of working}) definita all'interno delle Norme di Progetto
    \item Mantenere un canale di comunicazione costante, trasparente e proattivo sia con l'azienda proponente sia con il committente
\end{itemize}

\section{Analisi dei ruoli}
\subsection{Responsabile}
Il Responsabile ha il compito di governare il team, pianificare le risorse orarie ed economiche e coordinare le attività per assicurare il rispetto delle tempistiche predefinite e delle aspettative del committente. Rappresenta il progetto verso l'esterno, accentrando su di sé le responsabilità di scelta e approvazione finale.

\subsection{Amministratore di Sistema}
L'Amministratore ha il compito di definire, predisporre e mantenere l'ambiente di lavoro IT e le risorse informatiche necessarie a supporto del \textit{way of working} del team. Si occupa della gestione proattiva dell'infrastruttura, del tracciamento e della risoluzione dei ticket di disservizio e del controllo sulle procedure stabilite nelle Norme di Progetto. Attraverso l'implementazione di automatismi e l'ottimizzazione dei processi operativi, garantisce l'efficienza del gruppo di lavoro.

\subsection{Analista}
L'Analista è presente principalmente nelle fasi iniziali del ciclo di vita del software, si occupa della raccolta, dell'elaborazione e della formalizzazione dei requisiti attraverso un continuo confronto con il committente.

Il suo compito essenziale è delineare con chiarezza i requisiti di sistema, i vincoli e gli standard di accettazione che guideranno le successive attività di progettazione e sviluppo.

Tale concentrazione temporale nelle prime settimane del progetto comporta un monte ore ridotto ma di fondamentale rilevanza strategica.

\subsection{Progettista}
Il Progettista ha il compito di delineare la struttura del sistema affinché risponda ai requisiti identificati, selezionare le tecnologie implementative e fornire una documentazione esaustiva che guidi l'attività dei programmatori.

L'elevata rilevanza strategica e la complessità tecnica della progettazione comportano un monte ore consistente per questo ruolo nel corso del progetto.

\subsection{Programmatore}
Il Programmatore è la figura operativa preposta alla realizzazione concreta del prodotto software.

Costituendo la categoria più numerosa all'interno del team, ha il compito di tradurre in codice eseguibile la struttura architetturale e i componenti delineati dal Progettista. Pur possedendo specifiche competenze tecniche, agisce con deleghe decisionali circoscritte, uniformandosi alle scelte tecnologiche e alle convenzioni definite nelle Norme di Progetto.

Poiché la fase di scrittura e integrazione del codice risulta la più impegnativa a livello temporale, a questo ruolo viene allocato il monte ore più consistente dell'intero progetto.

\subsection{Verificatore}
Il Verificatore, dotato di un'approfondita conoscenza del \textit{way of working} del gruppo, possiede le capacità di giudizio e di relazione necessarie per effettuare controlli di qualità efficaci.

Si occupa di verificare la correttezza, l'esaustività e la coerenza di ogni documento redatto, oltre a testare e garantire la qualità del codice e del prodotto software durante la fase di codifica.

Poiché l'attività di controllo della qualità e conformità è continuativa, a questo ruolo viene assegnato un monte ore consistente.

\section{Criterio di rotazione dei ruoli}
La rotazione dei ruoli prevede che ogni persona assuma un solo ruolo per la durata di due settimane, con il cambio operativo che avviene formalmente ogni due lunedì. Il modello garantisce un'equa distribuzione del monte ore totale individuale, evitando disparità nel carico di lavoro e assicurando a tutti i componenti le medesime opportunità di apprendimento sui diversi ruoli. 

Al termine delle due settimane avviene il passaggio di consegne in cui il membro uscente trasferisce le informazioni e lo stato delle attività al subentrante.


\section{Preventivo dei Costi}

Di seguito è riportata la tabella riassuntiva dei costi preventivati per ciascun membro del Git Push Team e per ruolo, calcolata su un monte ore totale di 552 ore divise equamente (92 ore a testa). Segue il grafico della ripartizione oraria.

\vspace{0.5cm}

\begin{center}
    \resizebox{\textwidth}{!}{
    \renewcommand{\arraystretch}{1.3}
    \begin{tabular}{@{}l c c c c c c c@{}}
        \toprule
        \textbf{\color{gptinchiostro}Membro} & \textbf{\color{gptinchiostro}Resp.} & \textbf{\color{gptinchiostro}Amm.} & \textbf{\color{gptinchiostro}Analista} & \textbf{\color{gptinchiostro}Progett.} & \textbf{\color{gptinchiostro}Program.} & \textbf{\color{gptinchiostro}Verif.} & \textbf{\color{gptinchiostro}Totale} \\
        \midrule
        Enrico Biasini & 9 & 12 & 13 & 17 & 22 & 19 & 92 \\
        Luca Cesco Bolla & 9 & 12 & 13 & 17 & 22 & 19 & 92 \\
        Alessio Corrò & 9 & 12 & 13 & 17 & 22 & 19 & 92 \\
        Francesco Guglielmo & 9 & 12 & 13 & 17 & 22 & 19 & 92 \\
        Francesco Mazzurana & 9 & 12 & 13 & 17 & 22 & 19 & 92 \\
        Gaetano Spagnoli & 9 & 12 & 13 & 17 & 22 & 19 & 92 \\
        \midrule
        \textbf{Totale Ore} & \textbf{54} & \textbf{72} & \textbf{78} & \textbf{102} & \textbf{132} & \textbf{114} & \textbf{552} \\
        Costo Orario (\texteuro) & 30,00 & 20,00 & 25,00 & 25,00 & 15,00 & 15,00 & \\
        Costo Ruolo (\texteuro) & 1.620 & 1.440 & 1.950 & 2.550 & 1.980 & 1.710 & \\
        \bottomrule
    \end{tabular}
    }
\end{center}

\vspace{0.5cm}
\begin{center}
    \Large \textbf{\color{gptaccento}Costo Totale: \texteuro\ 11.250,00}
\end{center}

\vspace{1.5cm}

\newpage

\begin{center}
    \textbf{\color{gptinchiostro}Grafico degli impegni (Percentuale di ore di lavoro per ruolo)} \\[1.5em]
    \begin{tikzpicture}
        \def\radius{3.5cm}
        
        % Programmatore 23.91%
        \filldraw[fill=gptaccento!90, draw=black] (0,0) -- (0:\radius) arc (0:86.09:\radius) -- cycle;
        \node at (43.04:\radius*0.6) {23.91\%};
        \node at (43.04:\radius*1.3) {Programmatore};
        
        % Progettista 18.48%
        \filldraw[fill=gptaccento!75, draw=black] (0,0) -- (86.09:\radius) arc (86.09:152.61:\radius) -- cycle;
        \node at (119.35:\radius*0.6) {18.48\%};
        \node at (119.35:\radius*1.2) {Progettista};
        
        % Verificatore 20.65%
        \filldraw[fill=gptaccento!60, draw=black] (0,0) -- (152.61:\radius) arc (152.61:226.96:\radius) -- cycle;
        \node at (189.78:\radius*0.6) {20.65\%};
        \node at (189.78:\radius*1.3) {Verificatore};
        
        % Responsabile 9.78%
        \filldraw[fill=gptaccento!45, draw=black] (0,0) -- (226.96:\radius) arc (226.96:262.17:\radius) -- cycle;
        \node at (244.57:\radius*0.6) {9.78\%};
        \node at (244.57:\radius*1.2) {Responsabile};
        
        % Analista 14.13%
        \filldraw[fill=gptaccento!30, draw=black] (0,0) -- (262.17:\radius) arc (262.17:313.04:\radius) -- cycle;
        \node at (287.61:\radius*0.6) {14.13\%};
        \node at (287.61:\radius*1.1) {Analista};
        
        % Amministratore 13.04%
        \filldraw[fill=gptaccento!15, draw=black] (0,0) -- (313.04:\radius) arc (313.04:360:\radius) -- cycle;
        \node at (336.52:\radius*0.6) {13.04\%};
        \node at (336.52:\radius*1.4) {Amministratore};
        
    \end{tikzpicture}
\end{center}

\subsection{Motivazioni della divisione oraria}
Il preventivo orario complessivo di 552 ore è stato distribuito con l'obiettivo di ottimizzare la produttività e la qualità del prodotto finale, garantendo al contempo il contenimento dei costi. La logica alla base di questa ripartizione è la seguente:

\paragraph{Programmatore e Verificatore (132 e 114 ore):} Costituiscono il nucleo operativo del progetto. Si è scelto di allocare un monte ore sostanzioso anche alla verifica affinché il testing e la \textit{code review} procedano di pari passo con lo sviluppo, assicurando un alto standard qualitativo e arginando la proliferazione di bug.

\paragraph{Progettista (102 ore):} Un forte investimento orario è stato dedicato alla progettazione. Definire una base architetturale solida e ben strutturata è un passaggio critico per garantire l'estensibilità, l'affidabilità e la corretta integrazione tra i diversi moduli del sistema. Una progettazione chiara agevola il lavoro dei programmatori e abbatte i rischi di rifattorizzazione.

\paragraph{Analista (78 ore):} Il lavoro di analisi si concentra per lo più nelle fasi esplorative e di redazione dei casi d'uso. La quota ore garantisce uno studio attento del dominio aziendale richiesto dal capitolato.

\paragraph{Responsabile e Amministratore (54 e 72 ore):} L'apporto di gestione e amministrazione è stato contingentato per contenere il costo complessivo del progetto. Pur trattandosi di ruoli fondamentali, il loro impegno è focalizzato su interventi mirati e ad alta efficienza per il coordinamento delle attività, il supporto operativo e la gestione della configurazione.

\end{document}
